\textbf{Phần 1. Trắc nghiệm lựa chọn}
\setcounterpageref{ex}{0}

\begin{ex} %[2H7V35]
Tính chu vi hình tam giác có độ dài ba cạnh lần lượt là $18\text{ cm}, 25\text{ cm}, 27\text{ cm}$.
\choice
{$68\text{ cm}$}
{$69\text{ cm}$}
{\True $70\text{ cm}$}
{$71\text{ cm}$}
\loigiai{
$18 + 25 + 27 = 70\text{ (cm)}$.
\textbf{Đáp án của câu hỏi là C}
}
\end{ex}

\begin{ex} %[2H7V35]
Một thùng chứa $45\text{ L}$ dầu, người ta đã dùng hết $18\text{ L}$. Trong thùng còn lại bao nhiêu lít dầu?
\choice
{$26\text{ L}$}
{\True $27\text{ L}$}
{$28\text{ L}$}
{$25\text{ L}$}
\loigiai{
$45 - 18 = 27\text{ (L)}$.
\textbf{Đáp án của câu hỏi là B}
}
\end{ex}

\begin{ex} %[2H7V35]
Đoạn thẳng $AB$ dài $3\text{ dm } 5\text{ cm}$. Đổi độ dài đoạn thẳng $AB$ ra xăng-ti-mét là
\choice
{$305\text{ cm}$}
{$53\text{ cm}$}
{$30\text{ cm}$}
{\True $35\text{ cm}$}
\loigiai{
$3\text{ dm } 5\text{ cm} = 35\text{ cm}$.
\textbf{Đáp án của câu hỏi là D}
}
\end{ex}

\begin{ex} %[2H7V35]
Bao ngô nặng $42\text{ kg}$, bao gạo nặng hơn bao ngô $15\text{ kg}$. Bao gạo nặng bao nhiêu ki-lô-gam?
\choice
{\True $57\text{ kg}$}
{$55\text{ kg}$}
{$56\text{ kg}$}
{$58\text{ kg}$}
\loigiai{
$42 + 15 = 57\text{ (kg)}$.
\textbf{Đáp án của câu hỏi là A}
}
\end{ex}

\begin{ex} %[2H7V35]
Đường gấp khúc gồm $4$ đoạn thẳng bằng nhau, mỗi đoạn dài $8\text{ cm}$. Chiều dài đường gấp khúc là
\choice
{$28\text{ cm}$}
{$30\text{ cm}$}
{\True $32\text{ cm}$}
{$36\text{ cm}$}
\loigiai{
$8 \times 4 = 32\text{ (cm)}$.
\textbf{Đáp án của câu hỏi là C}
}
\end{ex}

\begin{ex} %[2H7V35]
Một chiếc đồng hồ chỉ $8$ giờ $15$ phút sáng. Sau $45$ phút nữa là mấy giờ?
\choice
{$8$ giờ $45$ phút}
{\True $9$ giờ sáng}
{$9$ giờ $15$ phút}
{$10$ giờ}
\loigiai{
$15$ phút $+ 45$ phút $= 60$ phút $= 1$ giờ. $8$ giờ $+ 1$ giờ $= 9$ giờ sáng.
\textbf{Đáp án của câu hỏi là B}
}
\end{ex}

\begin{ex} %[2H7V35]
Hình tứ giác có độ dài các cạnh là $15\text{ cm}, 18\text{ cm}, 22\text{ cm}, 25\text{ cm}$. Chu vi hình tứ giác là
\choice
{$78\text{ cm}$}
{$79\text{ cm}$}
{$82\text{ cm}$}
{\True $80\text{ cm}$}
\loigiai{
$(15 + 25) + (18 + 22) = 40 + 40 = 80\text{ (cm)}$.
\textbf{Đáp án của câu hỏi là D}
}
\end{ex}

\begin{ex} %[2H7V35]
Có $3$ can dầu loại $5\text{ L}$ và $1$ can loại $10\text{ L}$. Tổng số lít dầu là
\choice
{\True $25\text{ L}$}
{$20\text{ L}$}
{$30\text{ L}$}
{$15\text{ L}$}
\loigiai{
$5 \times 3 + 10 = 15 + 10 = 25\text{ (L)}$.
\textbf{Đáp án của câu hỏi là A}
}
\end{ex}

\begin{ex} %[2H7V35]
Con ngỗng nặng $6\text{ kg}$, con vịt nhẹ hơn con ngỗng $4\text{ kg}$. Cả hai con nặng bao nhiêu ki-lô-gam?
\choice
{$7\text{ kg}$}
{$9\text{ kg}$}
{\True $8\text{ kg}$}
{$10\text{ kg}$}
\loigiai{
Con vịt nặng: $6 - 4 = 2\text{ (kg)}$. Cả hai con nặng: $6 + 2 = 8\text{ (kg)}$.
\textbf{Đáp án của câu hỏi là C}
}
\end{ex}

\begin{ex} %[2H7V35]
Nếu ngày $1$ tháng $11$ là thứ Năm thì ngày $15$ tháng $11$ là thứ mấy?
\choice
{Thứ Tư}
{\True Thứ Năm}
{Thứ Sáu}
{Thứ Bảy}
\loigiai{
Từ ngày 1 đến ngày 15 là 14 ngày (đúng 2 tuần lễ). Do đó vẫn là thứ Năm.
\textbf{Đáp án của câu hỏi là B}
}
\end{ex}

\begin{ex} %[2H7V35]
Hình vuông có cạnh $20\text{ cm}$. Chu vi của hình vuông đó là
\choice
{$60\text{ cm}$}
{$70\text{ cm}$}
{$100\text{ cm}$}
{\True $80\text{ cm}$}
\loigiai{
$20 \times 4 = 80\text{ (cm)}$.
\textbf{Đáp án của câu hỏi là D}
}
\end{ex}

\begin{ex} %[2H7V35]
Một sợi dây dài $1\text{ m}$. Cắt bớt đi $35\text{ cm}$. Sợi dây còn lại dài bao nhiêu xăng-ti-mét?
\choice
{\True $65\text{ cm}$}
{$60\text{ cm}$}
{$75\text{ cm}$}
{$70\text{ cm}$}
\loigiai{
$1\text{ m} = 100\text{ cm}$. Còn lại: $100 - 35 = 65\text{ (cm)}$.
\textbf{Đáp án của câu hỏi là A}
}
\end{ex}

\textbf{Phần 2. Trắc nghiệm đúng sai}
\setcounterpageref{ex}{0}

\begin{ex} %[2H7V35]
Xét tính đúng sai của các nhận định về đại lượng và hình học:
\choiceTF
{\True Đổi 1 m = 10 dm = 100 cm.}
{1 kg bông nhẹ hơn 1 kg sắt.}
{\True Chu vi của một hình tam giác bằng tổng độ dài 3 cạnh của nó.}
{Một ngày có 12 giờ.}
\loigiai{
\begin{itemize}
    \item Ý a) Đúng: Quy đổi đơn vị đo độ dài chuẩn.
    \item Ý b) Sai: Cả hai đều nặng đúng 1 kg nên bằng nhau.
    \item Ý c) Đúng: Đúng theo định nghĩa.
    \item Ý d) Sai: Một ngày có 24 giờ.
\end{itemize}
\begin{itemize}
    \item a) \textbf{Đúng}.
    \item b) \textbf{Sai}.
    \item c) \textbf{Đúng}.
    \item d) \textbf{Sai}.
\end{itemize}
}
\end{ex}

\begin{ex} %[2H7V35]
Một bình nước chứa 18 L nước, người ta rót vào 3 can, mỗi can 4 L:
\choiceTF
{\True Số nước đã rót ra là 12 L.}
{\True Trong bình còn lại 6 L nước.}
{Số nước còn lại nhiều hơn số nước đã rót ra.}
{\True Lượng nước còn lại có thể rót đầy thêm 1 can 4 L và 1 can 2 L.}
\loigiai{
\begin{itemize}
    \item Ý a) Đúng: $4 \times 3 = 12\text{ L}$.
    \item Ý b) Đúng: $18 - 12 = 6\text{ L}$.
    \item Ý c) Sai: Còn 6 L, đã rót 12 L nên còn lại ít hơn.
    \item Ý d) Đúng: $4 + 2 = 6\text{ L}$. Đúng.
\end{itemize}
\begin{itemize}
    \item a) \textbf{Đúng}.
    \item b) \textbf{Đúng}.
    \item c) \textbf{Sai}.
    \item d) \textbf{Đúng}.
\end{itemize}
}
\end{ex}

\textbf{Phần 3. Trả lời ngắn}
\setcounterpageref{ex}{0}

\begin{ex} %[2H7V35]
Tính chu vi hình tam giác có độ dài các cạnh là 17 cm, 24 cm, 29 cm.
\shortans{70}
\loigiai{
\textbf{Đáp số:} $70$\\
$17 + 24 + 29 = 70\text{ cm}$.
}
\end{ex}

\begin{ex} %[2H7V35]
Một thùng chứa 50 L dầu, người ta múc ra 28 L dầu. Hỏi trong thùng còn lại bao nhiêu lít dầu?
\shortans{22}
\loigiai{
\textbf{Đáp số:} $22$\\
$50 - 28 = 22\text{ L}$.
}
\end{ex}

\begin{ex} %[2H7V35]
Một bao gạo nặng 45 kg, một bao ngô nặng 38 kg. Hỏi cả hai bao nặng bao nhiêu ki-lô-gam?
\shortans{83}
\loigiai{
\textbf{Đáp số:} $83$\\
$45 + 38 = 83\text{ kg}$.
}
\end{ex}

\begin{ex} %[2H7V35]
Đường gấp khúc gồm 4 đoạn thẳng có độ dài bằng nhau, mỗi đoạn dài 7 cm. Tính độ dài đường gấp khúc đó.
\shortans{28}
\loigiai{
\textbf{Đáp số:} $28$\\
$7 \times 4 = 28\text{ cm}$.
}
\end{ex}

\begin{ex} %[2H7V35]
Một buổi học bắt đầu lúc 7 giờ sáng và kết thúc lúc 11 giờ trưa. Hỏi buổi học kéo dài bao nhiêu giờ?
\shortans{4}
\loigiai{
\textbf{Đáp số:} $4$\\
$11 - 7 = 4$ giờ.
}
\end{ex}

\begin{ex} %[2H7V35]
Đoạn thẳng MN dài 1 m, đoạn thẳng PQ dài 48 cm. Hỏi đoạn thẳng MN dài hơn đoạn thẳng PQ bao nhiêu xăng-ti-mét?
\shortans{52}
\loigiai{
\textbf{Đáp số:} $52$\\
$100 - 48 = 52\text{ cm}$.
}
\end{ex}
